\documentclass[11pt,a4paper]{article}
\usepackage[margin=25mm,headheight=14pt]{geometry}
\usepackage[T1]{fontenc}
\usepackage[utf8]{inputenc}
\usepackage{lmodern,amsmath,amssymb}
\usepackage{xcolor,graphicx,booktabs,longtable,array}
\usepackage{microtype,enumitem,fancyhdr,titlesec}
\usepackage{hyperref}
\definecolor{accent}{HTML}{163B55}
\hypersetup{colorlinks=true,linkcolor=accent,urlcolor=accent,
 pdftitle={A certified lower bound for the ovals problem},pdfauthor={Matthew J. Colbrook},
 pdfsubject={AIM 1034: mathematical result and proof}}
\urlstyle{same}
\Urlmuskip=0mu plus 1mu
\setlength{\parindent}{0pt}
\setlength{\parskip}{0.35em}
\setlength{\emergencystretch}{3em}
\setlist{itemsep=0.25em,topsep=0.3em}
\providecommand{\tightlist}{\setlength{\itemsep}{0pt}\setlength{\parskip}{0pt}}
\titleformat{\section}{\normalfont\large\bfseries\color{accent}}{\thesection}{0.7em}{}
\titleformat{\subsection}{\normalfont\normalsize\bfseries}{\thesubsection}{0.6em}{}
\pagestyle{fancy}
\fancyhf{}
\fancyhead[L]{\small\sffamily AIM 1034}
\fancyhead[R]{\small\sffamily Matthew J. Colbrook}
\fancyfoot[C]{\thepage}
\renewcommand{\headrulewidth}{0.3pt}
\setcounter{secnumdepth}{0}
\allowdisplaybreaks[2]
\title{\textbf{A certified lower bound for the ovals problem}}
\author{Matthew J. Colbrook\\[3pt]
{\small Department of Applied Mathematics and Theoretical Physics}\\
{\small University of Cambridge, United Kingdom}\\
{\small\href{mailto:mjc249@maths.cam.ac.uk}{\texttt{mjc249@maths.cam.ac.uk}}}}
\date{3 November 2026}
\begin{document}
\maketitle
\begin{abstract}
We derive a universal lower bound greater than $0.96221$ for the lowest curvature Schrödinger eigenvalue of a smooth closed curve of length $2\pi$. Projection averaging and a variation constraint yield an explicit dual certificate with outward-rounded arithmetic. The conjectured sharp bound one in AIM 1034 remains open.
\end{abstract}
\noindent\textbf{Repository target.}
This manuscript concerns AIM problem 1034 at revision
\texttt{3fa50884d129}. Its scope is stated below.
\medskip

\hypertarget{1034-status}{%
\section{Status}\label{1034-status}}

This note does \textbf{not} establish the conjectured lower bound 1. It
gives a complete derivation of the explicit lower bound

\[
\lambda_1\!\left(-\frac{d^2}{ds^2}+\kappa(s)^2\right)
\ge C_0>0.96221.
\]

\[
C_0\in(0.962217776056501568384828,
        0.962217776056501568384870).
\]

The curve has length \(2\pi\); the bound holds for smooth closed curves
in \(\mathbb R^m\), \(m\ge2\), including the spatial curves in the
problem. The numerical enclosure is certified by
\texttt{verify\_ovals\_bound.py}, using only outward-rounded
standard-library Decimal arithmetic and explicit series remainders. The
proof relies on Denzler's established reduction of the global
variational problem to smooth strictly convex planar minimizers, and on
the projection-averaging inequality of Suragan. The latter is also
derived below for clarity. No claim of priority or independent peer
review is made.

The cited September 2026 bound is approximately 0.9618316523. The
additional input here is a total-variation constraint on the antipodal
arclength profile. The bound below remains strictly below 1, so problem
1034 stays unresolved.

\hypertarget{1034-geometric-reduction}{%
\section{1. Geometric reduction}\label{1034-geometric-reduction}}

Denzler proves that the infimum of the curvature Schrödinger eigenvalue
among closed curves of length \(2\pi\) is realized by a planar,
analytic, strictly convex curve with positive curvature; his relaxed
formulation also permits a digon of quotient 1. Consequently, a lower
bound below 1 proved for smooth strictly convex planar curves holds for
all the curves in the original problem. This is Theorem 1 and Theorems
6, 8, 9 in the arXiv version of his paper, with renumbering in the
published version.

We therefore consider an oriented strictly convex planar loop with
tangent

\[
T(s)=(\cos\phi(s),\sin\phi(s)),\qquad
\phi'(s)=\kappa(s)>0.
\]

Let \(s=\sigma(\theta)\) be the inverse tangent-angle coordinate and put
\(g(\theta)=\sigma'(\theta)\). Then

\[
g>0,\qquad \int_0^{2\pi}g(\theta)\,d\theta=2\pi,
\qquad \int_0^{2\pi} e^{i\theta}g(\theta)\,d\theta=0.
\tag{1}
\]

The last identity is closure of the original loop. Define the antipodal
arclength

\[
\ell(\alpha)=\sigma(\alpha+\pi)-\sigma(\alpha)
            =\int_\alpha^{\alpha+\pi}g(\theta)\,d\theta.
\tag{2}
\]

It obeys three useful constraints:

\[
\ell(\alpha)+\ell(\alpha+\pi)=2\pi,
\tag{3}
\]

\[
\int_0^\pi \ell(\beta-t)\sin t\,dt=2\pi,
\tag{4}
\]

and, on every interval \(I\) of length \(\pi\),

\[
\operatorname{TV}_{I}(\ell)
=\int_I|g(\alpha+\pi)-g(\alpha)|\,d\alpha
\le\int_I(g(\alpha+\pi)+g(\alpha))\,d\alpha
=2\pi.
\tag{5}
\]

For completeness, (4) follows by Fubini from the complex identity

\[
\int_0^\pi \ell(\beta-t)e^{it}\,dt=2\pi i.
\]

Indeed, the inner kernel after reversing the integrals is
\(i(1+e^{i(\beta-\theta)})\); the second term integrates to zero by (1).
The imaginary part is (4). Constraint (5) is the extra information used
in this note.

\hypertarget{1034-projection-averaging}{%
\section{2. Projection averaging}\label{1034-projection-averaging}}

For a real \(\psi\in H^1_{\rm per}(0,2\pi)\), set \(Y(s)=\psi(s)T(s)\).
Then

\[
|Y'(s)|^2=|\psi'(s)|^2+\kappa(s)^2\psi(s)^2.
\tag{6}
\]

For each angle \(\alpha\), the scalar projection
\(Y(s)\cdot(-\sin\alpha,\cos\alpha)\) vanishes at the two points with
tangent angles \(\alpha\) and \(\alpha+\pi\). Apply the sharp Dirichlet
interval Poincaré inequality separately on the two complementary arcs,
whose lengths are \(\ell(\alpha)\) and \(\ell(\alpha+\pi)\), and
integrate in \(\alpha\). The identity

\[
\int_0^{2\pi}|v\cdot(-\sin\alpha,\cos\alpha)|^2\,d\alpha
=\pi|v|^2
\]

gives

\[
\int_0^{2\pi}(|\psi'|^2+\kappa^2\psi^2)\,ds
\ge2\pi\int_0^{2\pi}J(\phi(s))\psi(s)^2\,ds,
\tag{7}
\]

where

\[
J(\beta)=\int_0^\pi
\frac{\sin^2t}{\ell(\beta-t)^2}\,dt.
\tag{8}
\]

To verify the factor in (7), one of the two arc contributions, after
reversing the order of integration, is
\(\pi^2\int J(\phi(s))\psi(s)^2\,ds\). The other is identical. Their sum
is divided by \(\pi\), the factor on the averaged left side. For complex
\(\psi\), apply the real inequality to its real and imaginary parts and
add.

\hypertarget{1034-symmetrizing-the-arclength-profile}{%
\section{3. Symmetrizing the arclength
profile}\label{1034-symmetrizing-the-arclength-profile}}

Fix \(\beta\), and write

\[
h(t)=\frac{\ell(\beta-t)}\pi,\qquad 0\le t\le\pi.
\]

Equations (3)--(5) give

\[
h>0,\quad h(0)+h(\pi)=2,\quad
\int_0^\pi h(t)\sin t\,dt=2,\quad
\int_0^\pi|h'(t)|\,dt\le2.
\tag{9}
\]

Set \(\widetilde h(t)=(h(t)+h(\pi-t))/2\). Convexity of
\(x\mapsto x^{-2}\) yields

\[
\int_0^\pi\frac{\sin^2t}{h(t)^2}\,dt
\ge\int_0^\pi\frac{\sin^2t}{\widetilde h(t)^2}\,dt.
\tag{10}
\]

The symmetrized function still has the same weighted integral and no
larger total variation. Let \(L=\pi/2\) and let \(v\) be its restriction
to \([0,L]\). By symmetry,

\[
v(0)=1,\qquad
\int_0^L v(t)\sin t\,dt=1,\qquad
\int_0^L|v'(t)|\,dt\le1.
\tag{11}
\]

The right side of (10) equals twice \[
F(v)=\int_0^L\frac{\sin^2t}{v(t)^2}\,dt.
\tag{12}
\]

\hypertarget{1034-a-rigorous-dual-certificate}{%
\section{4. A rigorous dual
certificate}\label{1034-a-rigorous-dual-certificate}}

All parameters in the following definition are exact:

\[
a=\frac{5467}{10000},\qquad
b=2a,\qquad c=\frac{54893}{50000},\qquad
\eta=\frac2{c^3}.
\tag{13}
\]

Notice \(0<a<b<c\). Define

\[
t_a=\arcsin((a/c)^3),\qquad
t_b=\arcsin((b/c)^3),
\]

and the positive comparison function

\[
v_*(t)=
\begin{cases}
a,&0\le t\le t_a,\\
c(\sin t)^{1/3},&t_a\le t\le t_b,\\
b,&t_b\le t\le L.
\end{cases}
\tag{14}
\]

This is a comparison function only; no assertion that it satisfies (11)
is needed. Put

\[
r(t)=\eta\sin t-\frac{2\sin^2t}{v_*(t)^3}.
\]

It is nonnegative on \([0,t_a]\), zero on \([t_a,t_b]\), and nonpositive
on \([t_b,L]\). Define

\[
A=\int_0^{t_a}r(t)\,dt,\qquad
B=-\int_{t_b}^{L}r(t)\,dt,
\qquad F_* = F(v_*).
\tag{15}
\]

Both \(A\) and \(B\) are nonnegative. Let \(q(t)=-\int_t^L r(s)\,ds\).
Thus \(q'=r\), \(q(L)=0\), \(q(0)=B-A\), and

\[
\|q\|_\infty=\max\{B,A-B\}.
\tag{16}
\]

For any \(v\) satisfying (11), integration by parts gives

\[
\begin{aligned}
\int_0^L r(t)v(t)\,dt
&=-q(0)-\int_0^Lq(t)v'(t)\,dt\\
&\ge A-B-\max\{B,A-B\}\\
&=-\max\{0,2B-A\}.
\end{aligned}
\tag{17}
\]

The supporting-line inequality for \(x\mapsto\sin^2t/x^2\) is

\[
\frac{\sin^2t}{v(t)^2}
\ge\frac{3\sin^2t}{v_*(t)^2}
    -\frac{2\sin^2t}{v_*(t)^3}v(t)
=\frac{3\sin^2t}{v_*(t)^2}
 +(r(t)-\eta\sin t)v(t).
\]

Integrate and use (11) and (17). This proves

\[
F(v)\ge D:=3F_*-\eta-\max\{0,2B-A\}.
\tag{18}
\]

It follows from (8), (10), and (12) that \(J(\beta)\ge2D/\pi^2\) for
every \(\beta\). Substitution into (7) proves the desired universal
bound with

\[
\boxed{C_0=\frac{4D}{\pi}.}
\tag{19}
\]

\hypertarget{1034-explicit-evaluation-and-numerical-certification}{%
\section{5. Explicit evaluation and numerical
certification}\label{1034-explicit-evaluation-and-numerical-certification}}

Let

\[
I_a=\frac{t_a}{2}-\frac{\sin(2t_a)}4,
\qquad
I_b=\frac{L-t_b}{2}+\frac{\sin(2t_b)}4,
\qquad
I_m=\int_{t_a}^{t_b}(\sin t)^{4/3}\,dt.
\]

Then the constants in (15) are

\[
F_* = \frac{I_a}{a^2}+\frac{I_m}{c^2}+\frac{I_b}{b^2},
\]

\[
A=\eta(1-\cos t_a)-\frac{2I_a}{a^3},
\qquad
B=\frac{2I_b}{b^3}-\eta\cos t_b.
\tag{20}
\]

The certificate encloses these values as follows (digits below are
abbreviated, not the full outward intervals):

\begin{longtable}[]{@{}lr@{}}
\toprule\noalign{}
Quantity & Approximate value \\
\midrule\noalign{}
\endhead
\bottomrule\noalign{}
\endlastfoot
\(A\) & 0.00384988849391648837 \\
\(B\) & 0.00192548312815568169 \\
\(D\) & 0.75572407410315354459 \\
\(C_0\) & 0.96221777605650156838 \\
\end{longtable}

To avoid dependence on an uncertified special-function routine, put
\(r_a=a/c\), \(r_b=b/c\), and \(d_j=\binom{2j}{j}/4^j\le1\). The two
needed expansions are

\[
\arcsin(r^3)=\sum_{j\ge0}\frac{d_j r^{6j+3}}{2j+1},
\]

\[
I_m=\sum_{j\ge0}\frac{3d_j}{6j+7}
       (r_b^{6j+7}-r_a^{6j+7}).
\tag{21}
\]

After \(n\) terms, valid upper bounds for the positive tails are,
respectively,

\[
\frac{r^{6n+3}}{(2n+1)(1-r^6)},
\qquad
\frac{3r_b^{6n+7}}{(6n+7)(1-r_b^6)}.
\tag{22}
\]

The program uses \(n=2000\) and 65-digit outward rounding. It encloses
\(\pi\) by the alternating Machin formula and bounds the square roots
\(\cos t_a=\sqrt{1-r_a^6}\), \(\cos t_b=\sqrt{1-r_b^6}\) using the
documented correctly rounded Decimal square-root operation with a
deliberately enlarged error interval. It obtains

\[
0.9622177760565015683848281814
< C_0 <
0.9622177760565015683848697651.
\tag{23}
\]

In particular \(C_0>0.96221\) is a rigorous numerical comparison.

\hypertarget{1034-scope-and-the-remaining-gap}{%
\section{6. Scope and the remaining
gap}\label{1034-scope-and-the-remaining-gap}}

The projection-averaging estimate (7) still loses information about how
close all scalar projections are to their first Dirichlet interval
modes. The certificate improves only the pointwise lower bound on \(J\).
It therefore gives no justification for replacing \(C_0\) by 1.

Numerically optimizing the analogous constrained scalar problem suggests
a best constant around 0.9622183981, only slightly above the rational
certificate used here. That numerical observation is not needed in the
proof and is not asserted as a certified optimum. It suggests that a
substantial further step must retain more than this pointwise profile
information.

\hypertarget{1034-references}{%
\section{References}\label{1034-references}}

\begin{enumerate}
\def\labelenumi{\arabic{enumi}.}
\tightlist
\item
  J. Denzler, \emph{Existence and regularity for a curvature dependent
  variational problem}, Transactions of the AMS 367 (2015), 3829--3845.
  \href{https://arxiv.org/abs/1301.6322}{arXiv:1301.6322}.
\item
  D. Suragan, \emph{An Improved Bound for the Ovals Problem},
  arXiv:2609.10775v1, September 2026. Proposition 3.1 supplies (7),
  Lemma 2.2 supplies (4), and Section 4 discusses the reduction to
  arbitrary spatial loops.
  \href{https://arxiv.org/abs/2609.10775}{Primary manuscript}.
\end{enumerate}
\end{document}
